\textbf{Phần 1. Trắc nghiệm lựa chọn}
\setcounterpageref{ex}{0}

\begin{ex} %[4H2V08]
Cho một hình tam giác có một góc là góc tù. Hai góc còn lại của hình tam giác đó chắc chắn là
\choice
{Một góc vuông và một góc nhọn}
{Hai góc tù}
{\True Hai góc nhọn}
{Một góc nhọn và một góc tù}
\loigiai{
Trong một hình tam giác, tổng số đo của ba góc luôn bằng $180^\circ$.\\
Vì tam giác đã có một góc tù (lớn hơn $90^\circ$) nên tổng số đo của hai góc còn lại phải bé hơn $90^\circ$.\\
Do đó, cả hai góc còn lại bắt buộc đều phải là hai góc nhọn.\\
\textbf{Đáp án của câu hỏi là C}
}
\end{ex}

\begin{ex} %[4H2V08]
Trong các góc sau, góc nào có số đo lớn nhất?
\choice
{\True Góc bẹt}
{Góc tù bé nhất có số đo là số tự nhiên}
{Góc vuông}
{Góc nhọn lớn nhất có số đo là số tự nhiên}
\loigiai{
Ta có:\\
- Góc nhọn lớn nhất là $89^\circ$.\\
- Góc vuông là $90^\circ$.\\
- Góc tù bé nhất là $91^\circ$.\\
- Góc bẹt là $180^\circ$.\\
Vậy góc có số đo lớn nhất là góc bẹt.\\
\textbf{Đáp án của câu hỏi là A}
}
\end{ex}

\begin{ex} %[4H2V08]
Trong một hình thang cân có hai đáy không bằng nhau và hai cạnh bên không vuông góc với đáy, bốn góc của hình thang đó gồm
\choice
{4 góc nhọn}
{4 góc vuông}
{1 góc nhọn, 2 góc vuông, 1 góc tù}
{\True 2 góc nhọn và 2 góc tù}
\loigiai{
Hình thang cân (không phải hình chữ nhật) có hai góc kề đáy lớn bằng nhau và đều là góc nhọn, hai góc kề đáy nhỏ bằng nhau và đều là góc tù.\\
Vậy hình thang cân đó gồm 2 góc nhọn và 2 góc tù.\\
\textbf{Đáp án của câu hỏi là D}
}
\end{ex}

\begin{ex} %[4H2V08]
Lúc 3 giờ 15 phút, góc tạo bởi kim giờ và kim phút của đồng hồ là
\choice
{Góc vuông}
{\True Góc nhọn}
{Góc tù}
{Góc bẹt}
\loigiai{
Lúc 3 giờ 15 phút, kim phút chỉ đúng vào số 3.\\
Kim giờ đã di chuyển qua số 3 một đoạn nhỏ (tương ứng với $\frac{15}{60} \times 30^\circ = 7{,}5^\circ$).\\
Vì $0^\circ < 7{,}5^\circ < 90^\circ$ nên góc tạo bởi hai kim là góc nhọn.\\
\textbf{Đáp án của câu hỏi là B}
}
\end{ex}

\begin{ex} %[4H2V08]
Một góc tù có số đo lớn hơn góc vuông là $a$ độ và bé hơn góc bẹt là $b$ độ. Tổng $a + b$ bằng
\choice
{$180^\circ$}
{$60^\circ$}
{\True $90^\circ$}
{$45^\circ$}
\loigiai{
Gọi số đo của góc tù đó là $X$ ($90^\circ < X < 180^\circ$).\\
Theo đề bài ta có:
$$a = X - 90^\circ$$
$$b = 180^\circ - X$$
Cộng hai vế lại ta được:
$$a + b = (X - 90^\circ) + (180^\circ - X) = 180^\circ - 90^\circ = 90^\circ$$
\textbf{Đáp án của câu hỏi là C}
}
\end{ex}

\begin{ex} %[4H2V08]
Một hình tứ giác $ABCD$ có góc $\widehat{A} = 100^\circ$, góc $\widehat{B} = 110^\circ$, góc $\widehat{C} = 75^\circ$. Góc $\widehat{D}$ của tứ giác đó là
\choice
{\True Góc nhọn có số đo $75^\circ$}
{Góc tù có số đo $95^\circ$}
{Góc vuông}
{Góc nhọn có số đo $65^\circ$}
\loigiai{
Tổng bốn góc trong một tứ giác bằng $360^\circ$.\\
Số đo góc $\widehat{D}$ là:
$$\widehat{D} = 360^\circ - (100^\circ + 110^\circ + 75^\circ) = 360^\circ - 285^\circ = 75^\circ$$
Vì $75^\circ < 90^\circ$ nên góc $\widehat{D}$ là góc nhọn có số đo $75^\circ$.\\
\textbf{Đáp án của câu hỏi là A}
}
\end{ex}

\begin{ex} %[4H2V08]
Một tờ giấy hình vuông có 4 góc vuông. Người ta dùng kéo cắt thẳng cắt đi một góc của tờ giấy hình vuông đó (đường cắt không đi qua các đỉnh còn lại). Số góc vuông còn lại trên mảnh giấy lớn là
\choice
{4 góc vuông}
{5 góc vuông}
{2 góc vuông}
{\True 3 góc vuông}
\loigiai{
Tờ giấy hình vuông ban đầu có 4 góc vuông.\\
Khi cắt chéo đi 1 góc vuông thì tại vết cắt mới tạo thành một cạnh mới và 2 góc tù mới.\\
Do đó 3 góc vuông cũ không bị cắt vẫn giữ nguyên.\\
Vậy mảnh giấy lớn còn lại đúng 3 góc vuông.\\
\textbf{Đáp án của câu hỏi là D}
}
\end{ex}

\begin{ex} %[4H2V08]
Một góc nhọn có số đo là $x$ (độ). Nếu gấp đôi số đo của góc nhọn đó ta được một góc tù. Giá trị của $x$ có thể là
\choice
{$30^\circ$}
{\True $65^\circ$}
{$40^\circ$}
{$45^\circ$}
\loigiai{
Để góc ban đầu là góc nhọn thì $x < 90^\circ$.\\
Khi gấp đôi số đo ta được góc tù thì:
$$2 \times x > 90^\circ \implies x > 45^\circ$$
Do đó $45^\circ < x < 90^\circ$.\\
Trong các phương án đã cho, chỉ có $x = 65^\circ$ thỏa mãn ($2 \times 65^\circ = 130^\circ$ là góc tù).\\
\textbf{Đáp án của câu hỏi là B}
}
\end{ex}

\begin{ex} %[4H2V08]
Người ta ghép hai chiếc ê-ke cùng loại có các góc $90^\circ; 60^\circ; 30^\circ$ kề nhau sao cho hai góc $60^\circ$ ghép sát vào nhau. Góc tạo thành bởi hai cạnh ngoài cùng là
\choice
{Góc vuông}
{Góc nhọn}
{\True Góc tù}
{Góc bẹt}
\loigiai{
Số đo góc tạo thành bởi hai góc $60^\circ$ ghép sát nhau là:
$$60^\circ + 60^\circ = 120^\circ$$
Vì $90^\circ < 120^\circ < 180^\circ$ nên góc tạo thành là một góc tù.\\
\textbf{Đáp án của câu hỏi là C}
}
\end{ex}

\begin{ex} %[4H2V08]
Trong các hình dưới đây, hình nào có số lượng góc tù nhiều nhất?
\choice
{\True Hình lục giác đều}
{Hình ngũ giác đều}
{Hình thang có hai góc tù}
{Hình chữ nhật}
\loigiai{
Ta xét số góc tù của từng hình:\\
- Hình chữ nhật: 0 góc tù (có 4 góc vuông).\\
- Hình thang có hai góc tù: có 2 góc tù.\\
- Hình ngũ giác đều: có 5 góc tù (mỗi góc bằng $108^\circ$).\\
- Hình lục giác đều: có 6 góc tù (mỗi góc bằng $120^\circ$).\\
Vậy hình có số lượng góc tù nhiều nhất là hình lục giác đều.\\
\textbf{Đáp án của câu hỏi là A}
}
\end{ex}

\begin{ex} %[4H2V08]
Cho góc $\widehat{xOy}$. Kẻ thêm $2$ tia khác nằm giữa hai cạnh $Ox$ và $Oy$. Hỏi có tất cả bao nhiêu góc đỉnh $O$ được tạo thành?
\choice
{$3$ góc}
{$4$ góc}
{$5$ góc}
{\True $6$ góc}
\loigiai{
Kẻ thêm 2 tia nằm giữa hai tia $Ox$ và $Oy$ thì tại đỉnh $O$ có tất cả:
$$2 + 2 = 4\text{ (tia chung gốc)}$$
Số góc tạo thành là:
$$4 \times (4 - 1) : 2 = 6\text{ (góc)}$$
\textbf{Đáp án của câu hỏi là D}
}
\end{ex}

\begin{ex} %[4H2V08]
Cho tam giác $ABC$ có góc đỉnh $A$ là góc tù. Khẳng định nào dưới đây là khẳng định đúng?
\choice
{Tam giác $ABC$ có thể có thêm một góc vuông}
{\True Hai góc đỉnh $B$ và $C$ đều là góc nhọn}
{Hai góc đỉnh $B$ và $C$ có tổng số đo lớn hơn $90^\circ$}
{Góc đỉnh $B$ có thể là góc tù}
\loigiai{
Tổng ba góc trong tam giác bằng $180^\circ$.\\
Vì góc đỉnh $A$ là góc tù ($> 90^\circ$) nên tổng hai góc đỉnh $B$ và $C$ phải nhỏ hơn $90^\circ$.\\
Do đó cả hai góc đỉnh $B$ và $C$ đều phải là các góc nhọn.\\
\textbf{Đáp án của câu hỏi là B}
}
\end{ex}

\textbf{Phần 2. Trắc nghiệm đúng sai}
\setcounterpageref{ex}{0}

\begin{ex} %[4H2V08]
Xét tính đúng sai của các khẳng định sau về các loại góc:
\choiceTF
{\True Góc nhọn luôn có số đo bé hơn góc tù.}
{Tổng số đo của hai góc nhọn luôn là một góc nhọn.}
{\True Hiệu số đo giữa một góc bẹt và một góc vuông là một góc vuông.}
{\True Trong một hình tam giác không thể có hai góc tù.}
\loigiai{
\begin{itemize}
    \item Vì góc nhọn $< 90^\circ$ và góc tù $> 90^\circ$ nên góc nhọn luôn bé hơn góc tù. Khẳng định a Đúng.
    \item Ví dụ hai góc nhọn $60^\circ$ và $60^\circ$ có tổng bằng $120^\circ$ là góc tù. Khẳng định b Sai.
    \item Ta có: $180^\circ - 90^\circ = 90^\circ$ (góc vuông). Khẳng định c Đúng.
    \item Tổng ba góc tam giác bằng $180^\circ$, nếu có hai góc tù thì tổng hai góc đã lớn hơn $180^\circ$ (vô lí). Khẳng định d Đúng.
\end{itemize}
\begin{itemize}
    \item a) \textbf{Đúng}.
    \item b) \textbf{Sai}.
    \item c) \textbf{Đúng}.
    \item d) \textbf{Đúng}.
\end{itemize}
}
\end{ex}

\begin{ex} %[4H2V08]
Xét các góc tạo bởi kim giờ và kim phút trên mặt đồng hồ vào các thời điểm:
\choiceTF
{\True Lúc 1 giờ đúng, hai kim tạo thành một góc nhọn.}
{\True Lúc 7 giờ đúng, góc bé tạo bởi hai kim là một góc tù.}
{Lúc 3 giờ đúng, hai kim tạo thành một góc nhọn.}
{\True Lúc 6 giờ đúng, hai kim tạo thành một góc bẹt.}
\loigiai{
\begin{itemize}
    \item Lúc 1 giờ đúng: góc bằng $30^\circ$ (góc nhọn). Khẳng định a Đúng.
    \item Lúc 7 giờ đúng: góc bé gồm 5 khoảng giờ nên bằng $5 \times 30^\circ = 150^\circ$ (góc tù). Khẳng định b Đúng.
    \item Lúc 3 giờ đúng: góc bằng $90^\circ$ là góc vuông (không phải góc nhọn). Khẳng định c Sai.
    \item Lúc 6 giờ đúng: hai kim thẳng hàng tạo thành góc bẹt $180^\circ$. Khẳng định d Đúng.
\end{itemize}
\begin{itemize}
    \item a) \textbf{Đúng}.
    \item b) \textbf{Đúng}.
    \item c) \textbf{Sai}.
    \item d) \textbf{Đúng}.
\end{itemize}
}
\end{ex}

\textbf{Phần 3. Trả lời ngắn}
\setcounterpageref{ex}{0}

\begin{ex} %[4H2V08]
Cho $5$ tia chung gốc $O$. Hỏi có tất cả bao nhiêu góc đỉnh $O$ được tạo thành?
\shortans{10}
\loigiai{
\textbf{Đáp số:} $10$\\
Số góc tạo thành từ 5 tia chung gốc là:
$$5 \times (5 - 1) : 2 = 10\text{ (góc)}$$
}
\end{ex}

\begin{ex} %[4H2V08]
Một góc tù có số đo lớn gấp $3$ lần góc nhọn $40^\circ$. Hỏi số đo của góc tù đó bằng bao nhiêu độ? (Chỉ ghi số).
\shortans{120}
\loigiai{
\textbf{Đáp số:} $120$\\
Số đo của góc tù đó là:
$$40^\circ \times 3 = 120^\circ$$
}
\end{ex}

\begin{ex} %[4H2V08]
Hai góc kề nhau có tổng số đo bằng $180^\circ$ (tạo thành góc bẹt). Biết góc thứ nhất gấp $4$ lần góc thứ hai. Tính số đo của góc thứ nhất theo độ (chỉ ghi số).
\shortans{144}
\loigiai{
\textbf{Đáp số:} $144$\\
Tổng số phần bằng nhau là:
$$1 + 4 = 5\text{ (phần)}$$
Số đo của góc thứ nhất là:
$$180^\circ : 5 \times 4 = 144^\circ$$
}
\end{ex}

\begin{ex} %[4H2V08]
Cho hình tam giác $ABC$ có góc $\widehat{A} = 70^\circ$, góc $\widehat{B} = 50^\circ$. Tính số đo của góc đỉnh $C$ theo độ (chỉ ghi số).
\shortans{60}
\loigiai{
\textbf{Đáp số:} $60$\\
Tổng ba góc trong tam giác $ABC$ bằng $180^\circ$.\\
Số đo góc $\widehat{C}$ là:
$$\widehat{C} = 180^\circ - (70^\circ + 50^\circ) = 60^\circ$$
}
\end{ex}

\begin{ex} %[4H2V08]
Một hình tứ giác có ba góc nhọn lần lượt có số đo là $70^\circ; 80^\circ; 65^\circ$. Tính số đo của góc thứ tư theo độ (chỉ ghi số).
\shortans{145}
\loigiai{
\textbf{Đáp số:} $145$\\
Tổng bốn góc trong hình tứ giác bằng $360^\circ$.\\
Số đo của góc thứ tư là:
$$360^\circ - (70^\circ + 80^\circ + 65^\circ) = 145^\circ$$
}
\end{ex}

\begin{ex} %[4H2V08]
Một bạn học sinh vẽ một góc vuông đỉnh $O$, sau đó vẽ thêm $2$ tia nằm bên trong góc vuông đó. Hỏi có tất cả bao nhiêu góc nhọn đỉnh $O$ được tạo thành?
\shortans{5}
\loigiai{
\textbf{Đáp số:} $5$\\
Góc vuông ban đầu có 2 tia, vẽ thêm 2 tia bên trong nên có tất cả 4 tia chung gốc $O$.\\
Tổng số góc tạo thành từ 4 tia là:
$$4 \times (4 - 1) : 2 = 6\text{ (góc)}$$
Trong 6 góc này có 1 góc lớn nhất là góc vuông ban đầu ($90^\circ$).\\
Năm góc còn lại đều có số đo bé hơn $90^\circ$ nên đều là góc nhọn.\\
Vậy có tất cả 5 góc nhọn.
}
\end{ex}
